%-----------TEMPLATE INFO--------------
% Resume in Latex
% Author : Sourabh Bajaj
% Website: https://github.com/sb2nov/resume
% License : MIT
%------------------------
%-----------DOCUMENT INFO--------------
% Author : Marcus Handley
% Website:
%------------------------
\documentclass[a4paper,10pt]{article}

\usepackage{latexsym}
\usepackage[empty]{fullpage}
\usepackage{titlesec}
\usepackage{marvosym}
\usepackage[usenames,dvipsnames]{color}
\usepackage{verbatim}
\usepackage{enumitem}
\usepackage[pdftex]{hyperref}
\usepackage{fancyhdr}


\pagestyle{fancy}
\fancyhf{} % clear all header and footer fields
\fancyfoot{}
\renewcommand{\headrulewidth}{0pt}
\renewcommand{\footrulewidth}{0pt}

% Adjust margins
\addtolength{\oddsidemargin}{-0.375in}
\addtolength{\evensidemargin}{-0.375in}
\addtolength{\textwidth}{1in}
\addtolength{\topmargin}{-.4in}
\addtolength{\textheight}{1.0in}

\urlstyle{same}

\raggedbottom
\raggedright
\setlength{\tabcolsep}{0in}

\definecolor{green}{RGB}{0,145,147}

% Sections formatting
\titleformat{\section}{
  \vspace{-4pt}\scshape\raggedright\Large
}{}{0em}{\textcolor{green}}[\color{green}\titlerule \vspace{-5pt}]

%-------------------------
% Custom commands
\newcommand{\resumeItem}[2]{
  \item[--]
    \textbf{#1}{#2 \vspace{-2pt}}
  }


\newcommand{\resumeSubheading}[4]{
  \vspace{-1pt}\item
    \begin{tabular*}{0.97\textwidth}{l@{\extracolsep{\fill}}r}
      \textbf{#1} & \small{#2} \\
      \textbf{#3} &  \small{#4} \\
    \end{tabular*}\vspace{-5pt}
}

\newcommand{\resumeProjectHeading}[3]{
  \vspace{-1pt}\item
    \begin{tabular*}{0.97\textwidth}{l@{\extracolsep{\fill}}r}
      \textbf{#1} \\
      \small{#2} \\
      \small{\href{#3}{#3}} \\
    \end{tabular*}\vspace{-5pt}
}

\newcommand{\resumeSubItem}[2]{\resumeItem{\small{#1}}{#2}\vspace{-4pt}}

\renewcommand{\labelitemii}{ }

\newcommand{\resumeSubHeadingListStart}{\begin{itemize}[leftmargin=*]}
\newcommand{\resumeSubHeadingListEnd}{\end{itemize}\vspace{-0.4cm}}
\newcommand{\resumeItemListStart}{\begin{itemize}
    [labelsep=1ex, itemindent=0em, leftmargin=1em
    ]}

\newcommand{\resumeItemListEnd}{\end{itemize}\vspace{-5pt}}


%-------------------------------------------
%%%%%%  CV STARTS HERE  %%%%%%%%%%%%%%%%%%%%%%%%%%%%

\begin{document}

%----------HEADING-----------------
\begin{tabular*}{\textwidth}{l@{\extracolsep{\fill}}r}
	\textbf{\LARGE Luca Choteborsky} &
	\href{mailto:luca.choteborsky@outlook.com}{luca.choteborsky@outlook.com}\\
	Production Engineer at Meta | M.Eng Computer Science, University of Cambridge
	& +44 7889 291392 \\ \href{https://github.com/lucachot}{Github}
	\href{https://www.linkedin.com/in/luca-choteborsky-72ba70190/}{LinkedIn}
\end{tabular*}

\section{Profile}
Production Engineer at Meta; M.Eng Computer Science from the University of
Cambridge with coursework in ML Systems (PyTorch) and distributed data processing.
Built and operated distributed infrastructure supporting AI training clusters
at scale, with hands-on experience in Kubernetes orchestration.
\vspace{-0.2cm}

%-----------ACADEMICS------------
\section{Academics}
 {\renewcommand\labelitemi{}

  \resumeSubHeadingListStart

  \resumeSubheading
  {University of Cambridge}{October 2024 -- June 2025}
  {M.Eng Computer Science}{Distinction}
  \resumeItemListStart
  \resumeItem{}{Principles of Machine Learning Systems: Studied and
	  implemented algorithms and system techniques in PyTorch supporting
	  training and inference of ML models across a spectrum of computing systems.}
  \resumeItem{}{Large-scale data processing and optimisation: Hands-on
	  experience with dataflow programming using TensorFlow and cluster
	  computing; academic grounding in next-generation distributed systems.}
  \resumeItem{}{Dissertation: built and evaluated a federated Kubernetes
	  scheduler (Go) using online FPCA-based load prediction and distributed
	  gRPC aggregation, outperforming the default kube-scheduler on
	  utilisation balance and job completion time across ML inference workloads.}
  \resumeItemListEnd

  \resumeSubheading
  {University of Cambridge}{October 2021 -- June 2024}
  {BA Computer Science}{Double 1st Class}
  \resumeItemListStart
  \resumeItem{}{Machine Learning and Real-world Data: practicals in
	  statistical classification (Naive Bayes), sequence modelling (Hidden
	  Markov Models), and social network analysis; experimental methodology
	  including cross-validation, statistical testing, and evaluation design.}
  \resumeItem{}{Concurrent and Distributed Systems: concurrency primitives,
	  distributed algorithms, consistency models, and fault-tolerant
	  message-passing systems.}
  \resumeItem{}{Advanced Graphics and Image Processing: rendering equation
	  and path tracing, image-based rendering (light fields and Neural
	  Radiance Fields), optimization-based image processing, and colour
	  science; techniques directly underpinning modern computer vision
	  systems.}
  \resumeItemListEnd

  \resumeSubHeadingListEnd

  \vspace{-0.2cm}
  \section{Technical Skills}
  \resumeSubHeadingListStart
  \item
  \begin{tabular}{l@	{} l@{\hskip 1in}}
	  \textbf{Programming Languages} & \hspace{.5cm} Python, Go, C++, Rust, C, Java, SQL, OCaml, Bash, \LaTeX  \\
	  \textbf{ML \& Frameworks}      & \hspace{.5cm} PyTorch, TensorFlow, Kubernetes, Docker, gRPC
  \end{tabular}
  \resumeSubHeadingListEnd

  %-----------EXPERIENCE-----------------
  \vspace{-0.2cm}
  \section{Professional Experience}

   {\renewcommand\labelitemi{}
    \resumeSubHeadingListStart

    \resumeSubheading
    {Meta}{United Kingdom}
    {Production Engineer, Network AI PE}{September 2025 -- Present}
    \resumeItemListStart
    \resumeItem{NVLINK Bug Discovery \& Automated Repair: }{Root-caused a
	    configuration misapplication affecting \textgreater25\% of the NVLINK
	    switch fleet ($\sim$40,000 devices). Devised an automated mechanism to
	    detect impacted switches and apply targeted repair fleet-wide.}
    \resumeItem{Production Service Migration: }{Migrated APIs for a network
	    drain service handling \textgreater1 million daily requests; implemented
	    A/B testing to guarantee zero user traffic impact during transition.}
    \resumeItem{Drain Monitoring \& Observability: }{Used existing data
	    collection infrastructure to validate drain state across a fleet of
	    100,000+ switches; performed aggregations on collected data to prevent
	    alert noise from scaling with fleet size, minimising oncall burden.}
    \resumeItem{CEL Compiler CLI Tool (Go): }{Added CEL compilation to an
	    existing Go CLI tool, parsing CEL configuration files and generating
	    device domain filters. Enabled dev-time testing across all consuming
	    teams, increasing validation coverage without increasing canary
	    reservations.}
    \resumeItemListEnd

    \resumeSubheading
    {Zilliqa}{London, United Kingdom}
    {Intern, Software Engineering}{July -- September 2023}
    \resumeItemListStart
    \resumeItem{}{Modified the Ethereum block explorer Otterscan using React
	    and TypeScript; worked with Ethereum Virtual Machine APIs in C++ for
	    cross-chain interoperability.}
    \resumeItemListEnd

    \resumeSubHeadingListEnd
   }

  \vspace{-0.2cm}
  \section{Projects}
  {\renewcommand\labelitemi{}
   \resumeSubHeadingListStart

   \resumeProjectHeading
   {Pronto -- Federated Kubernetes Scheduler}
   {Go, Kubernetes Scheduling Framework, gRPC, Docker, FPCA}
   {https://github.com/lucachot/pronto-framework}
   \resumeItemListStart
   \resumeItem{}{Evaluated against the default kube-scheduler on CPU-intensive
	   and ML inference workloads, demonstrating improved node utilisation
	   balance and reduced job completion times under batch workloads.}
   \resumeItem{}{Built a distributed aggregation layer using gRPC streaming
	   where per-node FPCA estimates are continuously merged into a shared
	   global embedding --- a form of federated learning applied to
	   infrastructure load prediction.}
   \resumeItem{}{Implemented a custom Kubernetes scheduler plugin via the
	   official Scheduling Framework API, routing pods to nodes based on
	   real-time CPU and memory telemetry rather than declared resource requests.}
   \resumeItemListEnd

   \resumeProjectHeading
   {Light Field Viewer}
   {WebGL, GLSL, JavaScript}
   {https://github.com/lucachot/lfv-webgl}
   \resumeItemListStart
   \resumeItem{}{Built a WebGL/GLSL shader pipeline rendering interactive
	   light field scenes from multi-image arrays, with real-time aperture
	   control, focal plane adjustment, and XYZ/orbital camera navigation.}
   \resumeItem{}{Wrote a Python preprocessing pipeline to parse and reformat
	   raw light field image data (Stanford datasets) into the texture layout
	   consumed by the GPU renderer.}
   \resumeItemListEnd

   \resumeSubHeadingListEnd
  }

 }
\end{document}
